\documentclass[12pt]{article}
\usepackage[utf8]{inputenc}
\usepackage[T1]{fontenc}
\usepackage{hyperref}
\title{Wotu Language Resources for K-12 Curriculum Development}
\author{Compiled by Claude Sonnet 5.5 (claude-sonnet-5-5) for Andr\'as Kornai}
\date{}
\begin{document}
\maketitle

\section{Introduction}
Wotu (ISO 639-3 wtw; Glottocode wotu1240) is an endangered Austronesian language of the Wotu–Wolio branch of Celebic. It is spoken around Wotu town on the northern end of Bone Bay, in Luwu Timur regency, South Sulawesi, Indonesia. Search results give about 500 native speakers out of roughly 4,000 ethnic members (2009), while Joshua Project lists 6,000 people. The Luwu dialect of Bugis and Indonesian put heavy pressure on the language, and literacy in Wotu is limited. This compilation rests on a single round of web searches.

\section{Primary Educational Resources}
I found no primer, textbook or school programme for Wotu.

\section{Linguistic and Reference Materials}
\begin{itemize}
\item Overview: \url{https://en.wikipedia.org/wiki/Wotu_language}
\item An article from the York University Lobo journals site turned up in the results: \url{https://lobo.journals.yorku.ca/index.php/default/article/download/163/172/286}. I did not verify that it is about Wotu or what it contains, so it is not in the json.
\end{itemize}
I found no Wotu dictionary or reference grammar.

\section{Religious and Community Resources}
\begin{itemize}
\item Joshua Project entry for the Wotu people (it describes them as a Buginese subgroup, which conflicts with linguistic classification): \url{https://joshuaproject.net/people_groups/15419}
\end{itemize}

\section{Digital and Audio Resources}
\begin{itemize}
\item Endangered Languages Project entry: \url{https://endangeredlanguages.com/lang/2804}
\end{itemize}
No audio or software resources turned up.

\section{Curriculum Development Context}
Wotu is not known to be taught in schools. With so few speakers and Bugis and Indonesian dominant locally, any curriculum would need to begin with documentation and community-led revitalization.

\end{document}